\section{Left View and Right View of a Binary Tree}
\label{sec:trees-left-right-view}

\problemstatement{The right view is the sequence of nodes visible from the
right side of the tree: the \emph{last} node at each level (top to bottom).
The left view is the \emph{first} node at each level. Return both.}

\begin{patternbox}
\emph{``View from the side,'' ``first/last node per level''} is level-order
selection: for each level, the rightmost node is the right view, the
leftmost is the left view. This is even simpler than the HD-based views --
there is no horizontal distance, only ``which node ends (or begins) each
level.''
\end{patternbox}

\subsection*{Understanding the problem carefully}

Standing to the right of the tree, on each level you see only the node
furthest right, because it blocks the ones behind it. So the right view is
just ``the last node of each level.'' Equivalently, with a depth-first
approach that visits \emph{right child before left}, the first node seen
at each new depth is the right-view node -- which gives an elegant $O(h)$
recursion.

\begin{ideabox}[First Node Seen per Depth, Right-First DFS]
Recurse carrying the current depth. Visit the \emph{right} child before
the left. The first time the recursion reaches a given depth, the node it
is at is the rightmost at that depth -- record it. For the left view, do
the same but visit the \emph{left} child first.
\end{ideabox}

\subsection*{Algorithm (right view)}

\begin{enumerate}
  \item \textsc{DFS}(\textit{node}, \textit{depth}, \textit{out}):
  \begin{enumerate}
    \item If \textit{node} is null, return.
    \item If $\textit{depth} = \text{size of } \textit{out}$: this is the
          first node reached at this depth -- append its value.
    \item \textsc{DFS}(\textit{node}.right, \textit{depth}+1);
          \textsc{DFS}(\textit{node}.left, \textit{depth}+1).
  \end{enumerate}
\end{enumerate}

\subsection*{Dry run}

\dryrun{Tree: root $1$; left $2$ (right child $5$); right $3$.}

\begin{itemize}
  \item DFS($1$,$0$): out empty, depth $0=0 \Rightarrow$ record $1$.
        Recurse right into $3$, then left into $2$.
  \item DFS($3$,$1$): depth $1=$ size $1 \Rightarrow$ record $3$. (No
        children.)
  \item DFS($2$,$1$): depth $1 \ne$ size $2$ -- skip. Recurse right into
        $5$.
  \item DFS($5$,$2$): depth $2=$ size $2 \Rightarrow$ record $5$.
  \item Right view: $1,3,5$. (Left view by symmetry: $1,2,5$.)
\end{itemize}

\subsection*{C++ implementation}

\begin{lstlisting}
#include <bits/stdc++.h>
using namespace std;

struct TreeNode {
    int val;
    TreeNode* left;
    TreeNode* right;
    TreeNode(int v) : val(v), left(nullptr), right(nullptr) {}
};

void rightDFS(TreeNode* node, int depth, vector<int>& out) {
    if (node == nullptr) return;
    if (depth == (int)out.size()) out.push_back(node->val);
    rightDFS(node->right, depth + 1, out);   // right first
    rightDFS(node->left,  depth + 1, out);
}

void leftDFS(TreeNode* node, int depth, vector<int>& out) {
    if (node == nullptr) return;
    if (depth == (int)out.size()) out.push_back(node->val);
    leftDFS(node->left,  depth + 1, out);    // left first
    leftDFS(node->right, depth + 1, out);
}

int main() {
    TreeNode* root = new TreeNode(1);
    root->left = new TreeNode(2);
    root->right = new TreeNode(3);
    root->left->right = new TreeNode(5);

    vector<int> rv, lv;
    rightDFS(root, 0, rv);
    leftDFS(root, 0, lv);

    cout << "Right view: "; for (int x : rv) cout << x << " "; cout << "\n";
    cout << "Left view:  "; for (int x : lv) cout << x << " "; cout << "\n";
    return 0;
}
\end{lstlisting}

\begin{complexitybox}
\textbf{Time Complexity.} $O(n)$ -- every node is visited once.
\textbf{Space Complexity.} $O(h)$ for the recursion stack.
\end{complexitybox}

\begin{takeawaybox}
The ``record only when depth equals the current answer size'' test is a
compact way to detect ``first node reached at a new depth'' without an
explicit level loop. Swapping the child-visit order flips left view to
right view -- another example of one algorithm, two mirrored variants.
\end{takeawaybox}
